% !PRINTED 207
\weberpage{207}

1.~\gesperrt{To every system of conjugate poles of the group $G$
there corresponds an invariant form, whose degree
is equal to the number of poles of the system, and
whose roots are precisely these poles.}

We denote by $f_1, f_2, \ldots$ the invariant forms belonging
to the various systems of conjugate poles, which are of the
degrees $\mu, \mu', \ldots$ and which we shall call the \gesperrt{fundamental
forms} of the group.

If now $F(x_1, x_2)$ is any invariant form of the
group $G$ and $f(x, y)$ a fundamental form, and if $F(x, 1) = 0$
has a common root with $f(x, 1) = 0$, then all the roots
of $f = 0$ must at the same time be roots of $F = 0$. For
by hypothesis the two functions $F(x, 1)$ and
$F[\psi(x), 1]$ have the same roots. Hence if $F(a, 1) = 0$,
then also $F[\psi(a), 1] = F(a_1, 1) = 0$. We therefore have the
following theorem:

2.~\gesperrt{If an invariant form $F(x_1, x_2)$ belonging to $G$
has a divisor in common with one of the fundamental forms $f(x_1, x_2)$,
then $F(x_1, x_2)$ is divisible by $f(x_1, x_2)$.}

If $F(x_1, x_2)$ is an invariant form of the group $G$, and $\xi$
a root of the equation $F(x, 1) = 0$, then, as
$\psi$ runs through the elements of the group $G$, all the
$n$ quantities $\psi(\xi)$ are also roots of the same equation. If therefore the
degree of $F$ is lower than the degree of the group, these
quantities cannot all be different from one another, and it
follows for any two different substitutions
$\chi, \psi$ of $G$ that
\[
\chi(\xi) = \psi(\xi)
\]
or
\[
\xi = \chi^{-1}\psi(\xi),
\]
that is, $\xi$ must occur among the poles of the group $G$, and
$F$ is therefore, by Theorem 2., divisible by one of the fundamental
forms. Since we can apply the same mode of inference to the
quotient of the division, it follows:

3.~\gesperrt{An invariant form $F$ of the group $G$, whose
degree is lower than the degree of the group,
is, apart from a constant factor, a
product of powers of the fundamental forms.}
%CONT>>
